\subsection{Localizing a smooth family of minimizing enclosures}

Compact normal tests give an averaged trace identity. To turn it into
information about individual minimizing frontiers, we use their common
tangent plane at intersections. Positivity then prevents cancellation
between different frontiers. The remaining issue is contact with the
original boundary: the zero-expansion equation continues through contact,
and strong comparison identifies the boundary components that remain.

The next proposition takes the enclosure geometry as an explicit input.
Its regularity assumptions hold for the minimizing families considered
here in dimensions at most seven. They are not consequences merely of
smoothness of the original boundary in higher dimensions.

\begin{proposition}[Localization for smooth minimizing frontiers]
\label{prop:variation:smooth-localization}
Let $(\Omega,g,K)$ be smooth initial data on a connected manifold with
one end and nonempty compact smooth boundary
$S=\bigsqcup_{i=1}^r S_i$, where $r<\infty$ and each $S_i$ is connected.
Orient $S$ into $\Omega$, and suppose
\[
 H_S+\operatorname{tr}_S K=0,\qquad A=\operatorname{Area}_g(S)>0.
\]
A full enclosing cut means the entire compact boundary of a connected
closed exterior domain containing the sufficiently distant end, with
its manifold interior in $\operatorname{int}\Omega$. Its normal points
toward that end, and all components and all coincidence with $S$ count.
Assume $S$ is outer area-minimizing for these smooth cuts.

Let $\mathcal T$ be a measurable family of full minimizing frontiers,
equipped with a probability measure $\varpi$. Assume the following
geometric properties.
\begin{enumerate}
\item Every $T\in\mathcal T$ has area $A$, lies in one fixed compact
subset of $\Omega$, and is a compact embedded two-sided
$C^{1,\alpha}$ hypersurface without boundary, for some $0<\alpha<1$.
It is the entire boundary of a connected closed $C^1$ exterior domain
containing the distant end. Its manifold interior lies in
$\operatorname{int}\Omega$, and the coorientation of $T$ points toward
that end. In particular, whenever $T$ is smooth it is a full enclosing
cut of the stated comparison class.
\item The free part $T\setminus S$ is smooth and minimal. At contact
with $S$, its graph functions are locally $W^{2,2}$, its filled side
contains the original obstacle, and its normal agrees with that of $S$.
Thus its local graph lies on the exterior side of the obstacle graph.
\item The surface-area measures form a measurable kernel in $T$.
On the union of the free frontiers there is a single Borel tangent-plane
field $\Pi(x)$, continuous relative to that union, such that
$\Pi(x)=T_xT$ whenever $x\in T\setminus S$.
\end{enumerate}
Suppose also that
\begin{equation}
 \int_{\mathcal T}\int_T s\,\operatorname{tr}_{T_xT}K(x)
                   \,dA_g(x)\,d\varpi(T)=0
 \qquad\bigl(s\in C_c^\infty(\operatorname{int}\Omega)\bigr).
 \label{eq:variation:smooth-aggregate-trace}
\end{equation}
Then, for $\varpi$-almost every $T$, the whole frontier is a smooth
MOTS, its free part satisfies $H_T=\operatorname{tr}_T K=0$, and each
of its connected components is either an entire $S_i$ or is compactly
contained in $\operatorname{int}\Omega$.

Moreover, $\varpi\{T:T=S\}=1$ under either of the following hypotheses:
\begin{enumerate}
\item[(a)] $S$ is connected, and no compact smooth embedded two-sided
full enclosing cut lying wholly in $\operatorname{int}\Omega$, possibly
disconnected, has $\theta_+\le0$ everywhere.
\item[(b)] No full smooth enclosing cut has all components either equal
to entire original components $S_i$ or wholly in
$\operatorname{int}\Omega$, at least one component in that open exterior,
and $\theta_+\le0$ everywhere.
\end{enumerate}
In case \textup{(a)}, before outermostness is applied, almost every
frontier satisfies the sharper dichotomy: it equals $S$, or it is
wholly interior and satisfies $H_T=\operatorname{tr}_T K=0$ everywhere.
In either localization case, for every continuous symmetric tensor $h$,
\begin{equation}
 \int_{\mathcal T}\frac12\int_T\operatorname{tr}_{T_xT}h\,dA_g
                                  \,d\varpi(T)
       =\frac12\int_S\operatorname{tr}_S h\,dA_g.
 \label{eq:variation:smooth-localized-area}
\end{equation}
\end{proposition}

\begin{proof}
Define the aggregate area measure in the open exterior by
\[
 \beta(B)=\int_{\mathcal T}\operatorname{Area}_g(T\cap B)\,d\varpi(T)
 \qquad\bigl(B\subset\operatorname{int}\Omega\text{ Borel}\bigr).
\]
It is positive and has mass at most $A$. On the union of the free
frontiers put $F(x)=\operatorname{tr}_{\Pi(x)}K(x)$, and set $F=0$
elsewhere. This is a bounded measurable function on the common compact
region. Equation~\eqref{eq:variation:smooth-aggregate-trace} says that
the signed Radon measure $F\beta$ vanishes. Hence $F=0$
$\beta$-almost everywhere, and Tonelli's theorem gives
\[
 0=\int |F|\,d\beta
  =\int_{\mathcal T}\int_{T\setminus S}
          |\operatorname{tr}_{T_xT}K|\,dA_g\,d\varpi(T).
\]
For almost every $T$, its tangential trace therefore vanishes at
area-almost every free point. Each free part is smooth, so a nonzero
value would persist on an open patch of positive area. Thus its trace
vanishes everywhere there. Minimality gives $H_T=0$ on that part as
well. The common-plane hypothesis is essential: it makes the trace one
function of position before positivity is used.

Fix one of these frontiers and a contact point. Extend the smooth
one-sided coefficients of $g,K$ locally across $S$ only to write graph
equations. In a graph chart let $f$ describe $T$ and let the smooth
function $\psi$ describe the obstacle. Their ordered graphs have
matching normals at contact. On the coincidence set $Z=\{f=\psi\}$,
their first derivatives agree. Sobolev locality applied to
$\partial_j(f-\psi)\in W^{1,2}$ gives
$D^2f=D^2\psi$ almost everywhere on $Z$. Since $S$ is marginal, the
graph of $f$ satisfies $H+\operatorname{tr}_{\rm tan}K=0$ almost
everywhere on $Z$. It satisfies this equation on the free set by the
first part of the proof.

We verify that this almost-everywhere equation yields smoothness
through contact. In the chart it has the form
\[
 \partial_i\mathcal F^i(y,f,Df)=\mathcal B(y,f,Df),
\]
with smooth functions $\mathcal F,\mathcal B$. The matrix
$\mathcal F^i_{p_j}$ is uniformly elliptic on a smaller chart, since
$Df$ is bounded there. These are the usual graph coefficients for
mean curvature, with the smooth tangential trace of $K$ included in
$\mathcal B$. As $f\in W^{2,2}$, the divergence on the left is an
$L^2$ function. The almost-everywhere identity is therefore also a
weak divergence-form equation; there is no additional contact measure.
For $f_k=\partial_k f$, differentiation in distributions gives
\[
 \partial_i(a^{ij}\partial_j f_k)
   =\partial_i\bigl(\delta_{ik}\mathcal B
           -\mathcal F^i_{y_k}-\mathcal F^i_z f_k\bigr),
 \qquad a^{ij}=\mathcal F^i_{p_j}(y,f,Df).
\]
All displayed coefficients and the right-hand vector field are
$C^{0,\alpha}$, while $f_k\in W^{1,2}$. Interior divergence-form
Schauder regularity gives $f_k\in C^{1,\alpha}$ on smaller charts.
Thus $f\in C^{2,\alpha}$, and further differentiation and Schauder
estimates give smoothness; see~\cite{GilbargTrudinger2001}.
Consequently $T$ is a smooth MOTS across every contact point.

If $T$ touches $S_i$, the two ordered smooth graphs satisfy the same
zero-expansion equation with the same orientation. Their nonnegative
difference $w$ satisfies a linear uniformly elliptic equation
$a^{ij}w_{ij}+b^i w_i+cw=0$ with bounded coefficients. Replacing $c$
by $\min(c,0)$ gives an inequality with left side $\le0$, since
$w\ge0$. The strong minimum principle at a zero of $w$ forces local
coincidence. Therefore $T\cap S_i$ is both open and closed in the
connected hypersurface $S_i$, so $S_i\subset T$. Local coincidence
also makes $S_i$ open in $T$, and compactness makes it closed there;
it is one connected component of $T$. Every remaining component is
disjoint from $S$ and, by compactness, lies a positive distance inside
the open exterior. This proves the component assertion.

Suppose first that \textup{(a)} holds. If $T$ touches $S$, it contains
$S$ as a component. Since $\operatorname{Area}_g(T)=A=
\operatorname{Area}_g(S)$, no further component of positive area is
possible, and $T=S$. Otherwise $T$ is a wholly interior smooth full
enclosing MOTS, excluded by \textup{(a)}.

Under \textup{(b)}, any interior component of $T$ would give exactly
a prohibited partial-coincidence cut: all its other components are
interior or entire $S_i$, and $\theta_+=0$ everywhere. Hence
$T=\bigsqcup_{i\in I}S_i$ for some $I\subset\{1,\ldots,r\}$.
The area identity
\[
 \sum_{i\in I}\operatorname{Area}_g(S_i)
       =A=\sum_{i=1}^r\operatorname{Area}_g(S_i)
\]
and the positive area of every $S_i$ force $I=\{1,\ldots,r\}$.
Thus $T=S$ in this case too. All conclusions hold outside the single
null family discarded above. Integrating the area derivative of this
fixed frontier proves Equation~\eqref{eq:variation:smooth-localized-area}.
\end{proof}
